\newpage

\renewcommand\thesection{5}
\renewcommand\thesubsection{\arabic{subsection}}

\counterwithin{subsection}{section}
\setcounter{subsection}{0}

\begin{center}
    \section*{\textbf{BAB V \\ KESIMPULAN DAN SARAN}}
    \addcontentsline{toc}{section}{BAB V \quad KESIMPULAN DAN SARAN}
\end{center}

\subsection{Kesimpulan}
\hspace{0.5cm}Berdasarkan hasil kegiatan yang telah dilaksanakan, maka dapat ditarik kesimpulan bahwa agama memiliki peran penting dalam menanamkan moral antikorupsi kepada masyarakat. Hasil kuesioner menunjukkan bahwa mayoritas responden mengetahui bahwa agama mereka melarang korupsi secara eksplisit dan memahami adanya sanksi moral atau spiritual bagi pelaku korupsi. Hal ini menunjukkan bahwa pengetahuan keagamaan masyarakat mengenai larangan korupsi sudah cukup tinggi. Selain itu, hasil wawancara dengan tokoh agama Buddha, Kristen, Katolik, Hindu, dan Islam menunjukkan bahwa seluruh agama memandang korupsi sebagai pelanggaran moral yang serius karena berkaitan dengan tindakan mengambil hak orang lain secara tidak sah.

\hspace{0.5cm}Meskipun demikian, kegiatan ini juga menemukan adanya kesenjangan antara pengetahuan agama dan penerapan moral dalam kehidupan nyata. Masyarakat mengetahui bahwa korupsi merupakan tindakan yang salah, tetapi nilai tersebut belum selalu diwujudkan dalam perilaku sehari-hari. Hal ini terlihat dari skor mengenai dampak ceramah antikorupsi terhadap perilaku yang masih lebih rendah dibandingkan skor keyakinan bahwa korupsi bertentangan dengan nilai agama. Selain itu, pengetahuan responden mengenai tokoh agama yang aktif mengampanyekan antikorupsi juga masih belum terlalu tinggi. Dengan demikian, penanaman moral antikorupsi tidak cukup hanya dilakukan melalui penyampaian ajaran agama secara normatif, tetapi perlu diperkuat melalui pembinaan moral yang lebih kontekstual, keteladanan tokoh agama, serta pendidikan antikorupsi berbasis nilai agama.


\subsection{Saran}
\hspace{0.5cm} 
Berdasarkan hasil kegiatan yang telah dilakukan, terdapat beberapa saran yang dapat diberikan. Pertama, lembaga keagamaan diharapkan tidak hanya menyampaikan ajaran mengenai larangan korupsi secara umum, tetapi juga membahas kasus-kasus nyata yang dekat dengan kehidupan masyarakat, seperti suap, gratifikasi, penyalahgunaan wewenang, manipulasi data, mencontek, dan plagiarisme. Kedua, tokoh agama perlu lebih aktif hadir sebagai figur teladan dalam menyuarakan nilai kejujuran, tanggung jawab, dan integritas, baik di lingkungan keagamaan maupun di ruang publik. Ketiga, lembaga pendidikan, keluarga, dan masyarakat perlu bekerja sama dalam menanamkan pendidikan antikorupsi sejak dini agar nilai agama tidak hanya dipahami sebagai aturan, tetapi juga menjadi dasar dalam mengambil keputusan moral.

\hspace{0.5cm} 
Penelitian ini juga memiliki keterbatasan karena responden kuesioner belum sepenuhnya mewakili seluruh kelompok masyarakat secara merata. Selain itu, wawancara hanya dilakukan kepada lima tokoh agama sehingga masih terdapat kemungkinan untuk memperluas sudut pandang dengan melibatkan lebih banyak narasumber dari berbagai lembaga keagamaan, pendidikan, maupun masyarakat umum. Oleh karena itu, penelitian selanjutnya diharapkan dapat menggunakan jumlah responden yang lebih besar, latar belakang responden yang lebih beragam, serta pembahasan yang lebih mendalam mengenai pengaruh pendidikan agama terhadap pembentukan perilaku antikorupsi dalam kehidupan sehari-hari.
